\documentclass[tikz,border=10pt]{standalone}
\input{../../../docs/tikz/dag_preamble.tex}

\begin{document}
\begin{tikzpicture}[every node/.style={font=\sffamily\small}]

\dagPaperMetadata{Fundamental Limits of Testing the Independence of Irrelevant Alternatives in Discrete Choice}{Arjun Seshadri and Johan Ugander}{arXiv}{2020}{https://arxiv.org/abs/2001.07042v1}
\dagPaperLegendRightOfMetadata{
  \node[dag_model, dag_template_legend] (legModel) at (0,0) {formalized\\model or definition};
  \node[dag_lemma, dag_template_legend] (legLemma) at (4.7,0) {formalized\\lemma or Appendix support};
  \node[dag_result, dag_template_legend] (legResult) at (9.4,0) {formalized\\main result};
}
\daglegend{(legModel)(legLemma)(legResult)}{Legend}

\begin{dagPaperBody}
\begin{scope}[xshift=-0.2cm,yshift=-0.5cm]

\node[dag_model, text width=4.8cm] (Choice) at (0,0) {
  \textbf{Choice-system model}\\finite observed item/set incidences and joint probabilities
};
\node[dag_model, text width=4.8cm] (IIA) at (6.8,0) {
  \textbf{IIA model}\\Luce ratio representation and IIA projection
};
\node[dag_model, text width=4.8cm] (Testing) at (13.6,0) {
  \textbf{Separated testing model}\\total variation, finite product tests, minimax risk
};
\node[dag_model, text width=4.8cm] (Graph) at (20.4,0) {
  \textbf{Comparison-incidence model}\\bipartite graph, Eulerian balance
};
\node[dag_model, text width=4.8cm] (Cycles) at (27.2,0) {
  \textbf{Cycle-decomposition model}\\orientations, mean length, dispersion
};

\node[dag_model, text width=4.8cm] (Perturb) at (0,-4.6) {
  \textbf{Balanced perturbations}\\uniform null, signed alternatives, finite mixtures
};
\node[dag_lemma, text width=4.6cm] (Fact5) at (6.8,-4.6) {
  \textbf{Appendix Fact 5}\\convex hull of IIA is dense
};
\node[dag_lemma, text width=4.6cm] (Fact7) at (13.6,-4.6) {
  \textbf{Appendix Fact 7}\\entropy as minimum cross entropy
};
\node[dag_lemma, text width=4.6cm] (Lemma6) at (20.4,-4.6) {
  \textbf{Appendix Lemma 6}\\marginals determine the IIA KL projection
};
\node[dag_lemma, text width=4.6cm] (Lemma9) at (27.2,-4.6) {
  \textbf{Appendix Lemma 9}\\native bipartite short-cycle packing and residual bounds
};

\node[dag_lemma, text width=4.6cm] (Lemma2) at (0,-9.2) {
  \textbf{Lemma 2}\\separated orientations and no-harder test reduction
};
\node[dag_lemma, text width=4.6cm] (Lemma3) at (6.8,-9.2) {
  \textbf{Lemma 3}\\mixture chi-square bound
};
\node[dag_lemma, text width=4.6cm] (Lemma4) at (13.6,-9.2) {
  \textbf{Lemma 4}\\cycle-orientation chi-square exponent
};
\node[dag_lemma, text width=4.6cm] (Lemma10) at (27.2,-9.2) {
  \textbf{Appendix Lemma 10}\\two-tier Eulerian decomposition, mean and dispersion bounds
};

\node[dag_result, text width=5.2cm] (Theorem1) at (10.2,-14.2) {
  \textbf{Theorem 1}\\minimax-risk, testing-radius, and sample-complexity lower bounds
};
\node[dag_result, text width=5.0cm] (Cor1) at (20.4,-14.2) {
  \textbf{Corollary 1}\\global IIA-testing lower bound
};
\node[dag_lemma, text width=5.0cm] (Cor2) at (27.2,-14.2) {
  \textbf{Appendix Corollary 2}\\all-even-subsets logarithmic cycle bounds
};

\node[dag_template_note, text width=17.5cm] (ChuNote) at (13.6,-19.2) {
  \textbf{Chu et al. citation.} The source cites Chu et al. for standalone Appendix Lemma 8.
  No theorem from that citation is assumed or imported. The bipartite construction used by
  Appendix Lemma 10 and the corollaries is Appendix Lemma 9 and is proved natively in Lean.
};

\begin{scope}[on background layer]
\draw[dag_arrow] (Choice) -- (IIA);
\draw[dag_arrow] (Choice) -- (Perturb);
\draw[dag_arrow] (IIA.south west) -- (Fact5.north east);
\draw[dag_arrow] (Fact7) -- (Lemma6);
\draw[dag_arrow] (Choice.south east) -- ++(0.8,-1.0) -| (Lemma6.north west);
\draw[dag_arrow] (IIA.south) -- ++(0,-1.1) -| (Lemma6.north west);
\draw[dag_arrow] (Graph) -- (Lemma9);
\draw[dag_arrow] (Cycles.south) -- ++(0,-1.0) -| (Lemma9.north east);
\draw[dag_arrow] (Lemma9) -- (Lemma10);
\draw[dag_arrow] (Graph.south east) -- ++(1.0,-1.0) -| (Lemma10.north east);
\draw[dag_arrow] (Cycles) -- (Lemma10);
\draw[dag_arrow] (Perturb) -- (Lemma2);
\draw[dag_arrow] (Graph.south west) -- ++(-1.0,-1.4) -| (Lemma2.north east);
\draw[dag_arrow] (Cycles.south west) -- ++(-1.4,-1.8) -| (Lemma2.north east);
\draw[dag_arrow] (Perturb.south east) -- ++(0.9,-1.0) -| (Lemma3.north west);
\draw[dag_arrow] (Testing.south west) -- ++(-0.9,-1.0) -| (Lemma3.north east);
\draw[dag_arrow] (Lemma3) -- (Lemma4);
\draw[dag_arrow] (Cycles.south west) -- ++(-0.8,-2.0) -| (Lemma4.north east);
\draw[dag_arrow] (Lemma2.south) -- ++(0,-1.0) -| (Theorem1.north west);
\draw[dag_arrow] (Lemma3) -- (Theorem1);
\draw[dag_arrow] (Lemma4) -- (Theorem1);
\draw[dag_arrow] (Testing.south) -- ++(0,-1.3) -| (Theorem1.north east);
\draw[dag_arrow] (Theorem1) -- (Cor1);
\draw[dag_arrow] (Lemma10) -- (Cor1);
\draw[dag_arrow] (Lemma10) -- (Cor2);
\end{scope}

\end{scope}
\end{dagPaperBody}
\end{tikzpicture}
\end{document}
